\documentclass[11pt]{article}
\usepackage{amsmath,amssymb,amsthm}
\newtheorem{theorem}{Theorem}
\title{Path-centered original compact chains in growing scaled bath tubes}\author{}\date{}
\begin{document}\maketitle
\paragraph{Precise original conditional class.}
Use the kernels of the time dependent small bath transfer calculation, with $s(u)=1-\frac12\operatorname{ReTr}u$,
\[
 F_t(u,v)=\tfrac B2\sum_{r=1}^4s(uS_{t,r}^{-1})+
 Bs(ug_tv^{-1}h_t)+\tfrac B2\sum_{r=1}^4s(vS_{t+1,r}^{-1}),\qquad T_t=e^{-F_t}.
\]
Take a period $m\ge2$, arbitrary time-dependent bath, and a literal endpoint
conjugation $u_m=G u_0G^{-1}$. Endpoint onsite data must agree with this transport.
Every staple and twist has principal angle at most $\epsilon_B$, where
$\epsilon_B\to0$. In particular $\epsilon_B=P(B)/\sqrt B$ permits any growing
$P(B)=o(\sqrt B)$, including powers of $\log B$. Bath-only action remains outside
this fiber and is never dropped. Write $W=T_0\cdots T_{m-1}C_G$.

\begin{theorem}[Conditional product estimate]
For sufficiently small $\epsilon_B$ and sufficiently large $B$, uniformly in
$m$, the bath data, and $G$, the exact Perron eigenvalue $\rho(W)$ satisfies
\[
 \left|\operatorname{Tr}(W^q)/\rho(W)^q-1\right|
       \le C(2/5)^{mq-1},\qquad q\ge1.                                  \tag{1}
\]
The thresholds depend only on a fixed small angular neighborhood, not on a
bound for $\sqrt B\epsilon_B$. No claim that this tube is probable under the
original random bath is made.
\end{theorem}

\paragraph{1. A uniform location bound for every global minimizing path.}
Existence follows from compactness. At a site, with all other path variables
fixed, the energy depends on its unit quaternion through the inner product
with the sum of four onsite staple quaternions and two unit temporal-neighbor
quaternions. This remains true for $m=2$, including its rotated seam. (The
single-site self-edge case $m=1$ is deliberately excluded.) The onsite sum has
length at least $4\cos\epsilon_B$ and direction of angle at most $\epsilon_B$
from identity. Adding a vector of length at most two shows that every
coordinate of a minimizing path has angle at most
\[
 \epsilon_B+\arcsin\bigl(2/(4\cos\epsilon_B)\bigr)<\pi/4                  \tag{2}
\]
for sufficiently small $\epsilon_B$. Its coefficient is nonzero, and choosing
the direction opposite that sum would maximize rather than minimize the local
energy, so no second branch was admitted in(2).

Let $\theta$ be the largest path angle. Conjugation at the seam preserves this
angle. At a maximizing coordinate take the radial unit tangent pointing away
from identity. If $\theta\ge\epsilon_B$, each onsite staple's scalar product
with this tangent is at most $-\sin(\theta-\epsilon_B)$. A temporal neighbor,
after its two twists, has angle at most $\theta+2\epsilon_B$; its scalar product
with the same tangent is at most $\sin(2\epsilon_B)$. These assertions follow
by writing a quaternion as $(\cos\phi,\sin\phi\,n)$; the maximum tangent
projection is $\sin(\phi-\theta)$ and the angles in(2) keep this in its
monotone range. Stationarity therefore gives
\[
 4\sin(\theta-\epsilon_B)\le2\sin(2\epsilon_B),\qquad
 \theta\le\epsilon_B+\arcsin(\tfrac12\sin2\epsilon_B)\le3\epsilon_B.       \tag{3}
\]
All global minimizers lie in this same small product chart. In ordinary Lie
coordinates the Hessian of the unscaled total action at zero bath and zero
field is $6I-S-S^*\ge4I$, with orthogonal twisted shift $S$. Smoothness of the
local two-site functions and their bounded incidence imply a perturbation
bound $C\epsilon_B$ on this Hessian throughout the product ball of radius
$3\epsilon_B$ in each coordinate. Thus the total action is strictly convex
there for small $\epsilon_B$, uniformly in $m$, and the global minimum $u^*$
is unique. Only local finite-dimensional smooth functions are used in this
Hessian estimate; the operator norm bound does not sum over all sites.

\paragraph{2. Exact centering and telescoping the linear terms.}
Use $u=u_t^*\exp(i x\cdot\sigma/\sqrt B)$ at each time. These are left Haar
translations of exact principal coordinates, with
$c_B=(2\pi^2)^{-1}B^{-3/2}$ and
$j_B(x)=[\sin(|x|/\sqrt B)/(|x|/\sqrt B)]^2$ on $|x|<\pi\sqrt B$.
Set
\[
 f_t=F_t(u_t^*,u_{t+1}^*),\qquad
 p_t=\partial_x F_t(u_t^*e^{ix\cdot\sigma/\sqrt B},u_{t+1}^*)|_{x=0}.
\]
Path stationarity says that the second-variable derivative of $F_t$ is exactly
$-p_{t+1}$, with the literal rotation at the seam. Hence the centered scaled
kernel is
\[
 H_{B,t}(x,y)=\sqrt{j_B(x)j_B(y)}\,
 \exp[-F_t(u_t^*e^{ix\cdot\sigma/\sqrt B},u_{t+1}^*e^{iy\cdot\sigma/\sqrt B})
       +f_t+p_t\cdot x-p_{t+1}\cdot y].                                  \tag{4}
\]
The derivative subtraction in(4) is exact. The exponential multipliers at
successive sites cancel in the product. For each finite $B$ they are bounded
and invertible on the compact coordinate ball; their norms need not be
uniform. The translated period is similar to
$c_B^m e^{-\sum_t f_t}(H_{B,0}\cdots H_{B,m-1}C_G)$, including its endpoint
rotation. Thus the scalar is the actual minimum action, not its Taylor
approximation; powers multiply it exactly by $q$.

\paragraph{3. Uniform Gaussian domination after centering.}
For a fixed sufficiently small angular radius $\theta_0$, the Hessian of each
individual two-site $F_t$ in the scaled coordinates differs uniformly from
that of
$\frac32(|x|^2+|y|^2)-x\cdot y$ by $O(\theta_0+\epsilon_B)$ whenever
$|x|,|y|\le\theta_0\sqrt B$. The latter Hessian has minimum eigenvalue two.
Taylor's integral formula after the EXACT linear subtraction therefore bounds
the exponent's action below by $c(|x|^2+|y|^2)$ on this product ball.

For the complementary region, the unitary triangle inequality
$s(ab^{-1})\ge\frac12s(a)-s(b)$, applied to the four onsite terms and to the
center translation, gives a global lower bound
\[
 F_t-f_t-p_t\cdot x+p_{t+1}\cdot y
 \ge c_0(|x|^2+|y|^2)-C B\epsilon_B^2
                  -C\sqrt B\epsilon_B(|x|+|y|).                          \tag{5}
\]
Here $f_t\le CB\epsilon_B^2$ and $|p_t|\le C\sqrt B\epsilon_B$ follow from(3)
and smoothness. If either coordinate exceeds $\theta_0\sqrt B$, the two
negative terms in(5) are absorbed by its positive quadratic term by taking
$\epsilon_B/\theta_0$ small. Combining the two regions proves the global bound
\[
 0\le H_{B,t}(x,y)\le e^{-c(|x|^2+|y|^2)}.                             \tag{6}
\]
All constants are uniform in time, period length and allowed data.
On bounded $x,y$ sets, the Hessian in(4) tends uniformly to the flat Hessian,
because $\epsilon_B\to0$ and $B^{-1/2}\to0$. The canceled linear coefficients
may grow as $\sqrt B\epsilon_B$; they do not enter the remainder. Thus(6) and
bounded-set convergence imply uniform Hilbert--Schmidt/operator convergence
\[
 H_{B,t}\longrightarrow H_0(x,y)=e^{-\frac32(|x|^2+|y|^2)+x\cdot y}.       \tag{7}
\]
In particular this is not an $O(P(B))$ perturbation about one fixed origin.

\paragraph{4. Uniform trace norm, including coordinate-edge tails.}
A trace-norm bound is needed in addition to(7). Choose smooth cutoffs in each
coordinate equal to one for $|x|\le\theta_0\sqrt B/2$ and zero for
$|x|\ge\theta_0\sqrt B$. On this interior product, every derivative through
order four in the first variable of the cut off kernel is bounded by a fixed
polynomial in $|x|+|y|$ times the Gaussian in(6). This follows from the bounded
scaled Hessian and higher derivatives of the smooth local action; the exact
linear subtraction makes its first derivative grow at most linearly in
$|x|+|y|$. The Haar Jacobian is smooth with uniformly bounded derivatives on
this smaller angular ball. Hence, for the oscillator
$L_x=1-\Delta_x+|x|^2$, the kernel of $L_x^2H_{B,t}^{\rm interior}$ has
uniform $L^2$ norm. Since $L_x^{-2}$ is Hilbert--Schmidt in dimension three,
factorization gives a uniform trace norm for the interior operator.

The remaining terms have at least one endpoint outside the interior cutoff.
Use the ORIGINAL factorization into two onsite multipliers and the central
positive temporal convolution split as $Q^{1/2}UQ^{1/2}$. The $U$ here is the
left/right Haar twist, also after center translation. The diagonal of $Q$ is
one. Squared Hilbert--Schmidt norms are consequently onsite integrals with
one endpoint cut off. Outside the angular ball the onsite action is at least
$cB$, whereas all linear gauge multipliers are at most $e^{CB\epsilon_B}$
and the extracted scalar is at most $e^{CB\epsilon_B^2}$. An endpoint tail
therefore gives $e^{-c'B}$; the other factor has at worst
$e^{CB\epsilon_B+CB\epsilon_B^2}$. After dividing by $c_B$, the tail trace norm
still tends to zero, uniformly for $\epsilon_B\to0$. This avoids differentiating
the zero extension at the principal-coordinate boundary. Consequently
$\|H_{B,t}/\ell_0\|_1\le K$ uniformly.

\paragraph{5. Whole-product spectral conclusion and its exact scope.}
By(7), the normalized steps are within $1/20$ of $H_0/\ell_0$ for all allowed
large $B$. The real common-cone and exterior-square proof of the time dependent small bath transfer calculation now applies
VERBATIM to the centered kernels: cone aperture $1/4$, growth $15/16$,
Perron projection norm at most $17/15$, and exterior norm at most $121/400$.
The retained conjugation fixes the radial reference vector. It yields(1) for
the exact centered product eigenvalue. The similarity and scalar in(4) transfer
both its trace and eigenvalue back to the ORIGINAL $W$, with no accumulated
per-step error. No self-adjointness of the product was assumed.

This enlarges the conditional class from a fixed tiny scaled tube to any
vanishing angular tube, including logarithmically growing scaled radii.
It supplies neither a probability nor a relative boundary-overlap bound for
that class in the original Wilson bath. Integrating out these trajectories
still leaves their exact normalizers in the bath law. Global gauge/holonomy
sectors and the remaining coarse bath affinity cannot be discarded. Therefore
this is a conditional S94 ingredient, not an interacting S94 or mass-gap proof.
\end{document}
